% Source: Wald GR, physical PDF pp. 72--79
%         (printed pp. 59--66).
% Coverage: Section 4.2, Special Relativity; equations (4.2.1)--(4.2.39).
% Figures/tables/footnotes: Fig. 4.1; footnotes 2--4; no tables.
% Status: source-reviewed and compile-clean.
% Uncertainties: none.

\subsection{Special Relativity}\label{sec:4.2}

We have already described in section 1.2 the major revolution in our notions of
space and time brought about by the special theory of relativity. Here we will
reformulate these ideas in a manner closely analogous to the discussion of the
previous section.

In special relativity, it is assumed that spacetime has the manifold structure of
\(\mathbb R^4\). As already mentioned in section 1.2, it is assumed that there
exist preferred families of motion in spacetime, referred to as
\textquotedblleft inertial\textquotedblright{} or \textquotedblleft
nonaccelerating\textquotedblright{} motions. Inertial observers, it is further
assumed, can set up a rigid grid of metersticks, synchronize clocks placed at the
gridpoints, and label each event in spacetime by the grid coordinates
\(x^1,x^2,x^3\) and the clock reading \(t\) at an event. This map of
spacetime into \(\mathbb R^4\) is called a \emph{global inertial coordinate
system}. Many different global inertial coordinate systems are
possible---specifically, the different systems can be put into one-to-one
correspondence with elements of the ten-parameter Poincar\'e group. Thus, the
labels \(t=x^0,x^1,x^2,x^3\) of a given event do not have intrinsic
meaning. However, as already mentioned in section 1.3, the spacetime interval
\(I\) between two events \(x\) and \(\bar x\), defined (in units where
\(c=1\)) by
\begin{equation}
  I=-(x^0-\bar x^0)^2+(x^1-\bar x^1)^2
    +(x^2-\bar x^2)^2+(x^3-\bar x^3)^2 ,
  \tag{4.2.1}\label{eq:4.2.1}
\end{equation}
has the same value for all global inertial coordinate systems and thus can be
viewed as representing an intrinsic property of spacetime.

In parallel with our discussion of the previous section,
equation~\eqref{eq:4.2.1} suggests that we define the metric of spacetime
\(\eta_{ab}\) by
\begin{equation}
  \eta_{ab}=\sum_{\mu,\nu=0}^{3}
  \eta_{\mu\nu}(\dd x^\mu)_a(\dd x^\nu)_b
  \tag{4.2.2}\label{eq:4.2.2}
\end{equation}
with \(\eta_{\mu\nu}=\diag(-1,1,1,1)\), where \(\{x^\mu\}\) is any
global inertial coordinate system. Again, the tensor field \(\eta_{ab}\) thereby
obtained is independent of the choice of global inertial coordinates. Again, the
ordinary derivative operator, \(\partial_a\), of the global inertial coordinates
satisfies
\begin{equation}
  \partial_a\eta_{bc}=0
  \tag{4.2.3}\label{eq:4.2.3}
\end{equation}
and thus is the derivative operator associated with the spacetime metric. Again,
since ordinary derivatives commute, the curvature of \(\eta_{ab}\) vanishes. The
geodesics of \(\eta_{ab}\) are those curves which are straight lines when
expressed in global inertial coordinates. In particular, the timelike geodesics of
\(\eta_{ab}\) are precisely the world lines of inertial observers in spacetime.

Thus, the theory of special relativity asserts that \emph{spacetime is the manifold
\(\mathbb R^4\) with a flat metric of Lorentz signature defined on it}.
Conversely, the entire content of special relativity as we have presented it thus
far is contained in this statement, since, given \(\mathbb R^4\) with a flat
Lorentz metric, we can use the geodesics of this metric to construct global inertial
coordinates, etc.

Quantities of physical interest in special relativity again are represented by
tensor fields, but now quantities which in prerelativity physics were viewed as
spatial tensors are recognized in special relativity to comprise spacetime tensors.
In the context of special relativity, the principle of general covariance states that
the spacetime metric, \(\eta_{ab}\), is the only quantity pertaining to spacetime
structure which can appear in any physical law. This principle is believed to apply
to the laws of physics in special relativity with one important modification.
Experiments demonstrating parity violation and (indirectly) demonstrating the
failure of time reversal symmetry have shown that two further aspects of spacetime
structure can appear in physical laws: the time orientation and space orientation
of spacetime. By the \textquotedblleft time orientation\textquotedblright{} we
mean a continuous choice throughout spacetime of which half of the light cone
represents the future direction, and which half represents the past direction (see
Fig. 1.2, and see the beginning of chapter 8 for further discussion of this notion
in the context of curved spacetime). By the \textquotedblleft space
orientation\textquotedblright{} we mean a continuous choice of
\textquotedblleft right handed\textquotedblright{} versus \textquotedblleft left
handed\textquotedblright{} orthonormal triads of spacelike vectors at each point,
or equivalently, the specification of a continuous, nonvanishing totally
antisymmetric tensor field \(e_{abc}=e_{[abc]}\) on spacetime satisfying
\(t^ae_{abc}=0\) for some timelike vector field\footnote{Note that if \(t^a\)
is chosen to be continuous, thus defining a time orientation, then
\(\epsilon_{abcd}=-t_{[a}e_{bcd]}\) yields a spacetime orientation in the
sense of appendix B. Any two of (i) a time orientation, (ii) a space orientation,
and (iii) a spacetime orientation of Minkowski spacetime determines the third.}
\(t^a\). Thus, these two additional aspects of spacetime structure apparently do
enter the laws of physics. Similarly, the laws of physics in special relativity are
believed to satisfy the principle of special covariance with respect to the
\emph{proper} Poincar\'e transformations, i.e., with respect to the translations,
rotations, and boosts of spacetime, but not with respect to time reflections, parity
transformations, or their composition. Although these latter
\textquotedblleft improper\textquotedblright{} transformations are isometries,
they do not preserve all the relevant structure of spacetime since they reverse the
time and/or space orientation of spacetime.

Let us describe more explicitly the formulation of some of the laws of physics in
special relativity. We have already mentioned that curves are classified as
timelike, null, or spacelike according to whether the norm
\(\eta_{ab}T^aT^b\) of their tangent is, respectively, negative, zero, or
positive. Special relativity asserts that the paths in spacetime of material
particles are always timelike curves. (This fact is, of course, simply a
reformulation of the familiar statement that \textquotedblleft nothing can travel
faster than the speed of light\textquotedblright.) We may parameterize timelike
curves by the \emph{proper time} \(\tau\) defined by
\begin{equation}
  \tau=\int\bigl(-\eta_{ab}T^aT^b\bigr)^{1/2}\,\dd t ,
  \tag{4.2.4}\label{eq:4.2.4}
\end{equation}
where \(t\) is an arbitrary parameterization of the curve (with increasing \(t\)
corresponding to \textquotedblleft forward in time\textquotedblright), and
\(T^a\) is the tangent to the curve in this parameterization. According to
special relativity, \(\tau\) is precisely the time which would elapse on a clock
carried along the given curve. Different timelike curves connecting the same pair
of events may have different elapsed times (the \textquotedblleft twin
paradox\textquotedblright), just as different paths between two points of space can
have different lengths. As discussed in section 3.3, the maximum elapsed time
between two events is given by geodesic (i.e., inertial) motion.

The tangent vector \(u^a\) to a timelike curve parameterized by \(\tau\) is
called the \emph{4-velocity} of the curve. It follows directly from the definition
of \(\tau\) that the 4-velocity has unit length,
\begin{equation}
  u^au_a=-1 .
  \tag{4.2.5}\label{eq:4.2.5}
\end{equation}
As already mentioned above, a particle subject to no external forces will travel
on a geodesic; i.e., its 4-velocity will satisfy the equation of motion
\begin{equation}
  u^a\partial_a u^b=0 ,
  \tag{4.2.6}\label{eq:4.2.6}
\end{equation}
where \(\partial_a\) is the derivative operator associated with \(\eta_{ab}\),
i.e., the ordinary derivative operator of a global inertial coordinate system. If
forces are present, the right-hand side of equation~\eqref{eq:4.2.6} will be
nonzero (see, e.g., equation~\eqref{eq:4.2.26} below).

All material particles have an attribute known as \textquotedblleft rest
mass\textquotedblright{} \(m\), which appears as a parameter in the equations of
motion when forces are present. The energy-momentum 4-vector, \(p^a\), of a
particle of mass \(m\) is defined by
\begin{equation}
  p^a=mu^a .
  \tag{4.2.7}\label{eq:4.2.7}
\end{equation}
The \emph{energy} of a particle as measured by an observer---present at the site
of the particle---whose 4-velocity is \(v^a\) is defined by
\begin{equation}
  E=-p_av^a .
  \tag{4.2.8}\label{eq:4.2.8}
\end{equation}
Thus, in special relativity, energy is recognized to be the
\textquotedblleft time component\textquotedblright{} of the 4-vector \(p^a\).
For a particle at rest with respect to the observer (i.e., \(v^a=u^a\)),
equation~\eqref{eq:4.2.8} reduces to the familiar formula \(E=mc^2\) (in
our units with \(c=1\)). Since the spacetime metric, \(\eta_{ab}\), is flat
and thus parallel transport is path independent, we may define the energy of a
particle as measured by an observer who is not present at the site of the particle
to be the energy measured by the observer who is at the site of the particle and
has 4-velocity parallel to that of the distant observer.

Continuous matter distributions in special relativity are described by a symmetric
tensor \(T_{ab}\) called the \emph{stress-energy-momentum tensor}. For an
observer with 4-velocity \(v^a\), the component \(T_{ab}v^av^b\) is
interpreted as the energy density, i.e., the mass-energy per unit volume, as
measured by this observer. For normal matter, this quantity will be nonnegative,
\begin{equation}
  T_{ab}v^av^b\geq 0 .
  \tag{4.2.9}\label{eq:4.2.9}
\end{equation}
If \(x^a\) is orthogonal to \(v^a\), the component \(-T_{ab}v^ax^b\) is
interpreted as the momentum density of the matter in the \(x^a\)-direction. If
\(y^a\) also is orthogonal to \(v^a\), then \(T_{ab}x^ay^b\) represents
the \(x^a\)-\(y^a\) component of the stress tensor of the material defined near
the beginning of section 2.3 above. Thus, the stress tensor of prerelativity
physics is combined with the energy and momentum densities to form the
stress-energy-momentum tensor of special relativity. Following standard practice,
we often shall abbreviate the term \textquotedblleft stress-energy-momentum
tensor\textquotedblright{} as \textquotedblleft stress-energy
tensor\textquotedblright{} or even \textquotedblleft stress
tensor\textquotedblright.

A \emph{perfect fluid} is defined to be a continuous distribution of matter with
stress-energy tensor \(T_{ab}\) of the form
\begin{equation}
  T_{ab}=\rho u_au_b+P(\eta_{ab}+u_au_b) ,
  \tag{4.2.10}\label{eq:4.2.10}
\end{equation}
where \(u^a\) is a unit timelike vector field representing the 4-velocity of the
fluid. According to the above interpretation of \(T_{ab}\), the functions
\(\rho\) and \(P\) are, respectively, the mass-energy density and pressure of
the fluid as measured in its rest frame. The fluid is called
\textquotedblleft perfect\textquotedblright{} because of the absence of heat
conduction terms and stress terms corresponding to viscosity.

The equation of motion of a perfect fluid subject to no external forces is simply
\begin{equation}
  \partial^aT_{ab}=0 .
  \tag{4.2.11}\label{eq:4.2.11}
\end{equation}
Writing out equation~\eqref{eq:4.2.11} in terms of \(\rho\), \(P\), and
\(u^a\), and projecting the resulting equation parallel and perpendicular to
\(u^b\), we find:
\begin{align}
  u^a\partial_a\rho+(\rho+P)\partial^au_a&=0 ,
    \tag{4.2.12}\label{eq:4.2.12}\\
  (P+\rho)u^a\partial_au_b+(\eta_{ab}+u_au_b)\partial^aP&=0 .
    \tag{4.2.13}\label{eq:4.2.13}
\end{align}
In the nonrelativistic limit, \(P\ll\rho\), \(u^\mu=(1,\mathbf{v})\), and
\(v\,\dd P/\dd t\ll\lvert\boldsymbol{\nabla} P\rvert\), these equations become
\begin{equation}
  \frac{\partial\rho}{\partial t}+\boldsymbol{\nabla}\mathbin{\cdot}(\rho\mathbf{v})=0 ,
  \tag{4.2.14}\label{eq:4.2.14}
\end{equation}
\begin{equation}
  \rho\left\{\frac{\partial\mathbf{v}}{\partial t}
  +(\mathbf{v}\mathbin{\cdot}\boldsymbol{\nabla})\mathbf{v}\right\}=-\boldsymbol{\nabla} P ,
  \tag{4.2.15}\label{eq:4.2.15}
\end{equation}
so we see that equation~\eqref{eq:4.2.11} reduces to the conservation of mass,
equation~\eqref{eq:4.2.14}, and Euler's equation~\eqref{eq:4.2.15}.

Equation~\eqref{eq:4.2.11} has an important physical interpretation. Consider a
family of inertial observers with parallel 4-velocities \(v^a\), so that
\(\partial_bv^a=0\). According to the above interpretation of \(T_{ab}\), the
quantity
\begin{equation}
  J_a=-T_{ab}v^b
  \tag{4.2.16}\label{eq:4.2.16}
\end{equation}
represents the mass-energy current density 4-vector of the fluid as measured by
these observers. Equation~\eqref{eq:4.2.11} implies
\begin{equation}
  \partial^aJ_a=0 .
  \tag{4.2.17}\label{eq:4.2.17}
\end{equation}
Using Gauss's law (see appendix B), equation~\eqref{eq:4.2.17} implies that
over the three-dimensional boundary, \(S\), of any four-dimensional spacetime
volume \(V\), we have
\begin{equation}
  \int_S J_an^a\,\dd S=0 ,
  \tag{4.2.18}\label{eq:4.2.18}
\end{equation}
where \(n^a\) is the unit normal whose direction is defined in appendix B.
Applying this to the volume shown in Figure~\ref{fig:4.1}, we see that the energy
change in the fluid in this volume (i.e., the contribution to the integral,
equation~\eqref{eq:4.2.18}, from the top and bottom parts of \(S\)) equals
the time integrated energy flux into the volume (i.e., the contribution from the
\textquotedblleft side\textquotedblright{} part of \(S\)). Thus,
equation~\eqref{eq:4.2.18} implies conservation of energy. Conversely,
conservation of energy as measured by all inertial observers requires
equation~\eqref{eq:4.2.11}. Thus, more generally, conservation of energy implies
that equation~\eqref{eq:4.2.11} must hold for all continuous matter
distributions, not just for perfect fluids.

% Source: Wald GR, physical PDF p. 76 (printed p. 63).
% Coverage: Figure 4.1, spacetime volume and mass-energy current density.
% Figures/tables/footnotes: Figure only; no tables or footnotes.
% Status: source-reviewed and compile-clean.
% Uncertainties: none.
\begingroup
\renewcommand{\thefigure}{4.1}
\begin{figure}[htbp]
  \centering
  \begin{tikzpicture}[>=Stealth,scale=0.88]
    % Reference axes.
    \draw[->] (-4.35,2.4) -- (-4.35,3.7) node[above] {\(\text{time}\)};
    \draw[->] (-4.35,2.4) -- (-3.15,2.4) node[right] {\(\text{space}\)};
    \draw[->] (-4.35,2.4) -- (-4.90,1.8);

    % Cylindrical spacetime volume and its boundary S.
    \draw (-1.8,2.7) arc[start angle=180,end angle=360,x radius=2.05,y radius=0.52];
    \draw (2.3,2.7) arc[start angle=0,end angle=180,x radius=2.05,y radius=0.52];
    \draw (-1.8,-1.7) arc[start angle=180,end angle=360,x radius=2.05,y radius=0.52];
    \draw[dashed] (2.3,-1.7) arc[start angle=0,end angle=180,x radius=2.05,y radius=0.52];
    \draw (-1.8,2.7) -- (-1.8,-1.7);
    \draw (2.3,2.7) -- (2.3,-1.7);

    % Representative mass-energy current vectors.
    \foreach \x/\y/\dx in {-1.2/-2.48/-.10,-.55/-2.55/.02,.15/-2.55/.08,
      .85/-2.52/.13,1.5/-2.45/.18,-1.2/-.9/-.12,-.45/-.8/.02,.25/-.85/.08,
      1.05/-.75/.13,1.7/-.65/.17,-1.1/.25/-.16,-.3/.15/.01,.5/.3/.08,
      1.25/.4/.15,-.8/1.25/-.12,.05/1.35/.02,.9/1.3/.1,1.65/1.45/.16} {
      \draw[->,thick] (\x,\y) -- ++(\dx,.72);
    }
    \draw[->,thick] (-1.9,-.35) -- ++(-.55,.65);
    \draw[->,thick] (2.35,.25) -- ++(.55,.7);
    \draw[->,thick] (-1.55,1.7) -- ++(-.45,.65);
    \draw[->,thick] (1.95,1.8) -- ++(.42,.7);
    \draw[->,thick] (-1.10,2.35) -- ++(-.18,.78);
    \draw[->,thick] (-.15,2.45) -- ++(-.02,.82);
    \draw[->,thick] (.82,2.43) -- ++(.15,.82);
    \draw[->,thick] (1.65,2.35) -- ++(.26,.75);

    \node[left] at (-2.65,.42) {\(S\)};
    \draw[->,thick,bend left=18] (-2.60,.34) to (-1.82,.18);
    \node at (.25,.45) {\(V\)};
    % Reserve trim clearance above the artwork when this float lands at page top.
    \path[draw=none] (0,4.8) circle[radius=.1pt];
  \end{tikzpicture}
  \caption{A spacetime diagram showing the mass-energy current density \(J^a\)
  in a volume \(V\) of spacetime.}
  \label{fig:4.1}
\end{figure}
\endgroup


To illustrate the description of fields in special relativity, we will give two
examples: the scalar field and the electromagnetic field. Although no classical
scalar field exists in nature, it is instructive for many purposes to consider a field
\(\phi\) satisfying the Klein-Gordon equation,
\begin{equation}
  \partial^a\partial_a\phi-m^2\phi=0 .
  \tag{4.2.19}\label{eq:4.2.19}
\end{equation}
The stress-energy tensor of this scalar field is\footnote{General prescriptions for
determining the stress-energy tensor of a field are discussed below, in appendix E.}
\begin{equation}
  T_{ab}=\partial_a\phi\,\partial_b\phi
  -\frac12\eta_{ab}\left(\partial^c\phi\,\partial_c\phi+m^2\phi^2\right) .
  \tag{4.2.20}\label{eq:4.2.20}
\end{equation}
Again, \(T_{ab}\) satisfies the energy condition,
equation~\eqref{eq:4.2.9}, and is conserved,
equation~\eqref{eq:4.2.11}, by virtue of the field
equation~\eqref{eq:4.2.19}.

In prerelativity physics, the electric field \(\mathbf{E}\) and magnetic field
\(\mathbf{B}\) each are spatial vectors. In special relativity these fields are
combined into a single spacetime tensor field \(F_{ab}\) which is antisymmetric
in its indices, \(F_{ab}=-F_{ba}\). Thus \(F_{ab}\) has six independent
components. For an observer moving with 4-velocity \(v^a\), the quantity
\begin{equation}
  E_a=F_{ab}v^b
  \tag{4.2.21}\label{eq:4.2.21}
\end{equation}
is interpreted as the electric field measured by that observer, while
\begin{equation}
  B_a=-\frac12\epsilon_{ab}{}^{cd}F_{cd}v^b
  \tag{4.2.22}\label{eq:4.2.22}
\end{equation}
is interpreted as the magnetic field, where \(\epsilon_{abcd}\) is the totally
antisymmetric tensor of positive orientation with norm
\(\epsilon_{abcd}\epsilon^{abcd}=-24\) (see appendix B) so that in a
right-handed orthonormal basis we have \(\epsilon_{0123}=1\).

In terms of \(F_{ab}\), Maxwell's equations take the simple and elegant form,
\begin{align}
  \partial^aF_{ab}&=-4\pi j_b ,
    \tag{4.2.23}\label{eq:4.2.23}\\
  \partial_{[a}F_{bc]}&=0 ,
    \tag{4.2.24}\label{eq:4.2.24}
\end{align}
where \(j^a\) is the current density 4-vector of electric charge. Note that the
antisymmetry of \(F_{ab}\) implies that
\begin{equation}
  0=\partial^b\partial^aF_{ab}=-4\pi\partial^bj_b .
  \tag{4.2.25}\label{eq:4.2.25}
\end{equation}
Thus, Maxwell's equations imply \(\partial^bj_b=0\), which, by the same
argument as given above for \(J_a\), states that electric charge is conserved. The
equation of motion of a particle of charge \(q\) moving in the electromagnetic
field \(F_{ab}\) is
\begin{equation}
  u^a\partial_au^b=\frac{q}{m}F^b{}_cu^c ,
  \tag{4.2.26}\label{eq:4.2.26}
\end{equation}
which reformulates the usual Lorentz force law in terms of \(F_{ab}\).

The stress-energy tensor of the electromagnetic field is
\begin{equation}
  T_{ab}=\frac{1}{4\pi}\left\{
  F_{ac}F_b{}^c-\frac14\eta_{ab}F_{\dd e}F^{\dd e}\right\} .
  \tag{4.2.27}\label{eq:4.2.27}
\end{equation}
Again, \(T_{ab}\) satisfies the energy condition,
equation~\eqref{eq:4.2.9}, and if \(j^a=0\), we have
\(\partial^aT_{ab}=0\) by virtue of Maxwell's equations. If \(j^a\neq0\),
then the stress-energy \(T_{ab}\) of the electromagnetic field alone is not
conserved, but the total stress-energy of the field and the charged matter is still
conserved.

By the converse of the Poincar\'e lemma (see appendix B),
equation~\eqref{eq:4.2.24} implies that there exists a vector field \(A^a\)
(called the \emph{vector potential}) such that
\begin{equation}
  F_{ab}=\partial_aA_b-\partial_bA_a .
  \tag{4.2.28}\label{eq:4.2.28}
\end{equation}
In terms of \(A^a\), Maxwell's equations become
\begin{equation}
  \partial^a(\partial_aA_b-\partial_bA_a)=-4\pi j_b .
  \tag{4.2.29}\label{eq:4.2.29}
\end{equation}
We have the gauge freedom of adding the gradient \(\partial_a\chi\) of a
function \(\chi\) to \(A_a\) since, by equation~\eqref{eq:4.2.28}, this
leaves \(F_{ab}\) unchanged. By solving the equation
\begin{equation}
  \partial^a\partial_a\chi=-\partial^bA_b
  \tag{4.2.30}\label{eq:4.2.30}
\end{equation}
for \(\chi\), we may make a gauge transformation to impose the Lorenz gauge
condition,
\begin{equation}
  \partial^aA_a=0 ,
  \tag{4.2.31}\label{eq:4.2.31}
\end{equation}
in which case, using the commutativity of derivatives in flat spacetime,
equation~\eqref{eq:4.2.29} becomes
\begin{equation}
  \partial^a\partial_aA_b=-4\pi j_b .
  \tag{4.2.32}\label{eq:4.2.32}
\end{equation}

We may seek solutions of Maxwell's equations of the form of a wave oscillating
with constant amplitude,
\begin{equation}
  A_a=C_a\exp(\ii S) ,
  \tag{4.2.33}\label{eq:4.2.33}
\end{equation}
where \(C_a\) is a constant vector field (i.e., constant norm and everywhere
parallel to itself) and the function \(S\) is called the \emph{phase} of the wave.
To yield a solution with \(j^a=0\), the phase must satisfy (from
equation~\eqref{eq:4.2.32})
\begin{align}
  \partial^a\partial_aS&=0 ,
    \tag{4.2.34}\label{eq:4.2.34}\\
  \partial_aS\,\partial^aS&=0 ,
    \tag{4.2.35}\label{eq:4.2.35}
\end{align}
and (from equation~\eqref{eq:4.2.31})
\begin{equation}
  C_a\partial^aS=0 .
  \tag{4.2.36}\label{eq:4.2.36}
\end{equation}
Now, for any function \(f\) on any manifold with a metric, the vector
\(\nabla_af\) is normal (i.e., orthogonal) to the surfaces of constant \(f\),
since for any vector \(t^a\) tangent to the surface we have
\(t^a\nabla_af=0\). Equation~\eqref{eq:4.2.35} states that the normal
\(k^a=\partial^aS\) to the surfaces of constant \(S\) is a null vector,
\(k_ak^a=0\). We call such a surface a \emph{null hypersurface}. Note that
null hypersurfaces satisfy the property that their normal vector is tangent to the
hypersurface. Differentiation of equation~\eqref{eq:4.2.35} yields
\begin{align}
  0&=\partial_b(\partial_aS\,\partial^aS)\notag\\
   &=2(\partial^aS)(\partial_b\partial_aS)\notag\\
   &=2(\partial^aS)(\partial_a\partial_bS)\notag\\
   &=2k^a\partial_ak_b ,
  \tag{4.2.37}\label{eq:4.2.37}
\end{align}
i.e., the integral curves of \(k^a\) are null geodesics. Indeed, since the
derivation of equation~\eqref{eq:4.2.37} from
equation~\eqref{eq:4.2.35} is valid in curved spacetime as well (with
\(\nabla_a\) replacing \(\partial_a\) everywhere), it follows that all null
hypersurfaces in a Lorentz spacetime are generated by null geodesics.\footnote{The
derivation of equation~\eqref{eq:4.2.37} is appropriate to the case in which
one has a function \(S\) for which all the hypersurfaces of constant \(S\) are
null. For a single null hypersurface, the vanishing of
\(\epsilon^{abcd}k_c\nabla_d(k_ek^e)\) shows that \(k^a\) is tangent to a
(non-affinely parameterized) null geodesic.} The \emph{frequency} of the wave
(i.e., minus the rate of change of the phase of the wave) as measured by an
observer with 4-velocity \(v^a\) is given by
\begin{equation}
  \omega=-v^a\partial_aS=-v^ak_a .
  \tag{4.2.38}\label{eq:4.2.38}
\end{equation}

The most important solutions of the form, equation~\eqref{eq:4.2.33}, are the
plane waves,
\begin{equation}
  S=\sum_{\mu=0}^{3}k_\mu x^\mu ,
  \tag{4.2.39}\label{eq:4.2.39}
\end{equation}
where \(\{x^\mu\}\) are global inertial coordinates and \(k_\mu\) are
constants (and thus \(k^a\) is a constant vector field). From the theory of
Fourier transforms, it follows that all well behaved solutions of Maxwell's
% SOURCE ERRATUM: printed ``spacial''; corrected to ``spatial.''
equations which go to zero sufficiently rapidly at large spatial distances can be
expressed as superpositions of plane waves.

Our analysis, above, suggests that electromagnetic disturbances, i.e., light
signals, propagate on null geodesics, equation~\eqref{eq:4.2.37}. This
supposition is indeed correct: The retarded Green's function for Maxwell's
equations has support confined to the past light cone (see, e.g., Jackson 1962), so
the electromagnetic radiation from a source which reaches event \(p\) depends
only on what the source does on the past light cone of \(p\). Thus, the
terminology of \textquotedblleft light cone\textquotedblright, which we have been
using all along, is justified by Maxwell's equations; light does indeed propagate
along the light cone.
